% frontmatter/como-usar.tex -- instrucciones para el adulto.
\section*{Cómo usar este cuaderno}

Este cuaderno es para una niña o un niño de siete u ocho años que ya lee
sola y ya suma y resta hasta 20. No enseña cuentas de mayores, sino las
ideas de la economía de cada día: que las cosas tienen valor, que el
dinero sirve para cambiar unas cosas por otras, que no se puede tener
todo y hay que elegir, que ahorrar es guardar hoy para mañana, que
alguien hace el pan que compramos y que las aceras y los parques los
pagamos entre todos. Es de la misma familia que \emph{Aprendo a leer} y
\emph{Aprendo los números}, y va con \emph{Leo con lupa}: los mismos
personajes (Lucía, su hermano Dani, Toby, la abuela Rosa, Sofía, Hugo)
y los mismos temas semana a semana.

Hay una página para cada día laborable del año, de lunes a viernes,
también en vacaciones. Diez minutos al día son suficientes. Cada página
sigue siempre la misma estructura:

\begin{itemize}[leftmargin=6mm]
\item \textbf{Fecha}, \textbf{Hecho} y \textbf{Notas}: para el adulto.
  Marcar cada día si lo ha hecho sola, con ayuda o con dificultad no
  cuesta nada, y al cabo del año es un registro de todo lo que ha
  avanzado.
\item \textbf{Hoy}, en el recuadro verde: la historia del día, dos o
  tres frases que lee la niña o el niño. Los lunes, debajo, la
  \textbf{palabra de la semana} y lo que quiere decir; el viernes se
  repasa, y al final del cuaderno están todas, en \emph{Mi diccionario
  de economía}.
\item Una \textbf{actividad}. Las instrucciones (en letra pequeña y
  gris) las lee un adulto, si hace falta.
\end{itemize}

Cada actividad está explicada, una a una, en la página siguiente.

\subsection*{Cómo ayudar}

{\small
\begin{itemize}[leftmargin=6mm, itemsep=1pt, topsep=2pt]
\item Hablad de la historia del día: ¿qué habrías hecho tú? En economía
  casi nunca hay una sola respuesta, y lo que se aprende es a pensar en
  lo que se gana y en lo que se deja.
\item Con monedas de verdad, las cuentas se entienden mejor: contad las
  que haya en casa, jugad a las tiendas, dejad que pague y cuente la
  vuelta cuando vayáis a comprar.
\item Equivocarse con poco dinero es la mejor manera de aprender a usar
  el mucho.
\item Al final de cada trimestre hay una medalla, y al final del
  cuaderno, un diploma y una \textbf{clave de respuestas} para las
  páginas que tienen una respuesta.
\end{itemize}}
\newpage
